% =============================================================================
% 9. EVALUATION PLAN (FORWARD-LOOKING)
% Target: 0.5 page (~250 words)
% Source: cap5 distilled to one short paragraph + forward reference
% Status: SKELETON
%
% IMPORTANT: This section is intentionally short. The full evaluation
% protocol and results are the subject of a companion empirical paper.
% Reviewer-deflection: it answers "what about evaluation?" without
% consuming the empirical material.
% =============================================================================

\section{Evaluation Plan}
\label{sec:evaluation-plan}

The empirical validation of PATHCAST is reported in a companion paper
(under review); we summarise the design here for completeness.

The protocol follows Design Science
Research~\cite{hevner2004, peffers2007} on a dataset of 48 open-source
repositories selected by stratified sampling across primary language,
project scale, governance model, and workflow maturity, totalling
approximately 227{,}000 commits. The experimental design comprises
two stages. \textbf{E1 (pipeline validation)} assesses event-log and
process-model quality (correlation coverage, activity coverage,
fitness, precision, generalisation). \textbf{E2 (forecasting
accuracy)} compares PATHCAST against four baselines --- historical
mean, throughput-based Monte Carlo~\cite{vacanti2015}, analytical
Markov, and historical empirical CDF --- using the continuous ranked
probability score (CRPS) as primary metric, mean and median absolute
error for point accuracy, the aggregate calibration error (ACE) for
distributional fidelity, and $85\%$ coverage for the operational
reporting interval. The ML component is evaluated with AUC-ROC, F1,
and Brier score across three feature sets (PM-only, MSR-only,
Combined). Statistical analysis uses Wilcoxon signed-rank tests with
Bonferroni correction
($\alpha_{\text{adj}} = 0.0125$) across repositories and Cliff's
$\delta$ for effect
size~\cite{kampenes2007systematic}. Temporal hold-out (70/30 split)
and leave-one-project-out within domain groups control for
information leakage and cross-repository generalisation.

This paper does not report empirical results. The reader is referred
to the companion paper for the full evaluation, including calibration
plots, CRPS box plots, and feature-importance rankings.
